\section{Reduction to an abelian surface with cyclic image}
\label{sec:reduction}

The main construction will produce a simple abelian surface inside a
smooth projective variety, with infinite cyclic image on fundamental
groups.  We first explain why that is enough.  This separates the
geometric obstruction to Steinness from the group-theoretic construction
that occupies the intervening sections.

\begin{proposition}\label{prop:large-reduction}
Let $A$ be a simple abelian surface and let $A\hookrightarrow S'$ be a
closed embedding in a smooth connected projective complex variety.
Suppose that
\[
 \im\bigl(\pi_1(A)\longrightarrow\pi_1(S')\bigr)\simeq\Z.
\]
There exist a smooth connected projective fourfold $X$ with large
fundamental group and a closed embedding $A\hookrightarrow X$ with
infinite cyclic image on fundamental groups, such that the universal
cover of $X$ is not Stein.
\end{proposition}

We begin with the equivalence between largeness and the condition on
the universal cover used in Theorem~\ref{thm:main}.

\begin{lemma}\label{red:compact-criterion}
Let $X$ be a smooth connected projective complex variety.  Its universal
cover contains no positive-dimensional compact complex-analytic
subvariety if and only if $X$ has large fundamental group.
\end{lemma}

\begin{proof}
Suppose that $Z\subset X$ has positive dimension and that
$\pi_1(Z^\nu)$ has finite image in $\pi_1(X)$.  The covering of $Z^\nu$
corresponding to the kernel has finite degree and is compact.  Its map
to $X$ lifts holomorphically to $\widetilde X$.  The image of the lift
is a positive-dimensional compact analytic subvariety, by the proper
mapping theorem.

Conversely, let $K\subset\widetilde X$ be a positive-dimensional
irreducible compact analytic subvariety.  Its image $Z$ in $X$ is
analytic by the same theorem and algebraic by Chow's theorem.  The
restriction $K\to Z$ is finite: it is proper and its fibers are discrete,
since the covering map is locally biholomorphic.  Thus
$K^\nu\to Z^\nu$ is finite and surjective.  Such a map has finite-index
image on fundamental groups.  To see the latter assertion, remove a
proper analytic subset of the target so that the map becomes a finite
unramified cover of smooth connected spaces.  On these open sets the
index is its covering degree.  Their fundamental groups surject onto
those of the normal spaces, by general position for loops; passing to
these quotients can only decrease the index.  But the composite
$\pi_1(K^\nu)\to\pi_1(X)$ is zero, because $K^\nu$ maps to
$\widetilde X$.  Hence $\pi_1(Z^\nu)$ has finite image in $\pi_1(X)$.
\end{proof}

The next elementary interpolation observation allows us to control all
fibers of a general complete intersection simultaneously.

\begin{lemma}\label{red:interpolation}
Let $W$ be a projective variety, $F\subset W$ a closed subscheme, and
$L$ a very ample line bundle.  Fix a positive integer $\ell$.  For all
sufficiently large $m$ and every set of distinct points
$p_1,\ldots,p_\ell\in W\setminus F$, the evaluation map
\[
 H^0(W,\mathcal I_F\otimes L^m)
 \longrightarrow \bigoplus_{i=1}^{\ell}L^m|_{p_i}
\]
is surjective.  The same bound on $m$ works for every such set.
\end{lemma}

\begin{proof}
Choose $m_0$ so that $\mathcal I_F\otimes L^{m_0}$ is globally
generated.  For each $i$, choose a section $a_i$ nonzero at $p_i$.
For every $j\ne i$, very ampleness gives a section of $L$ vanishing at
$p_j$ and nonzero at $p_i$.  Their product $b_i$ vanishes at all the
$p_j$ with $j\ne i$ and is nonzero at $p_i$.  The sections $a_ib_i$
therefore give a diagonal basis for evaluation in degree
$m_0+\ell-1$.  Multiplying each by a section of the remaining power
of $L$ nonzero at $p_i$ gives the assertion in every higher degree.
\end{proof}

\begin{proof}[Proof of Proposition~\ref{prop:large-reduction}]
Put $d=\dim S'$, choose an abelian variety $B$ of dimension $g=6$, and
write
\[
 W=S'\times B,\qquad F=A\times\{0\},\qquad
 c=d+g-4=d+2.
\]
Fix a very ample line bundle $L$ on $W$.  We shall take $X$ to be the
common zero locus of $c$ general sections of
$\mathcal I_F\otimes L^m$, for a sufficiently large $m$.
There are two requirements: $X$ must be smooth with the expected
fundamental group, and its fibers over $B$ must be finite away from $F$.

\smallskip
\noindent\emph{Smoothness and the fundamental group.}
For large $m$, the sections generate the ideal away from $F$ and their
normal derivatives generate
$N^*_{F/W}\otimes L^m|_F$.  Bertini's theorem gives smoothness away
from $F$.  Along $F$, smoothness is the independence of $c$ general
vectors in a bundle of rank $d+g-2=c+2$.  The rank-deficient matrices
have codimension $3$, larger than $\dim F=2$, so no rank loss occurs
for a general tuple.  The same count for each prefix of the tuple
shows that the successive intersections are smooth.  They are ample
hypersurfaces in their predecessors, with final dimension four.
The Lefschetz hyperplane theorem consequently gives connectedness and
an isomorphism, induced by inclusion,
\begin{equation}\label{red:fundamental-group}
 \pi_1(X)\xrightarrow{\ \sim\ }
 \pi_1(S'\times B)=\pi_1(S')\times\pi_1(B).
\end{equation}
Here and below we use sufficiently high powers of a very ample line
bundle to obtain both ideal generation and generation of the indicated
normal derivatives; these are the usual applications of Serre
vanishing and Bertini's theorem \cite{Hartshorne}.

\smallskip
\noindent\emph{Fibers away from the fixed surface.}
We claim that a general such tuple excludes four distinct points of
$W\setminus F$ in any one fiber of $W\to B$.  Excluding four points
is enough to exclude every positive-dimensional fiber component
outside $F$; the choice $g=6$ makes the following incidence count
strict.  The parameter space of
ordered quadruples in one fiber has dimension at most $4d+g=4d+6$.
By Lemma~\ref{red:interpolation}, requiring all $c$ equations to
vanish at such a quadruple imposes exactly $4c=4d+8$ independent
linear conditions on the space of tuples of sections.  The incidence
variety of tuples and quadruples thus has dimension strictly less than
the section parameter space.  Its constructible image has closure of
strictly smaller dimension as well.  A general tuple lies outside this
closure.  This condition is compatible with the open smoothness
conditions already imposed.

It follows that each fiber of $X\to B$ contains at most three points
outside $F$.  In particular, every positive-dimensional subvariety of
a fiber is contained in $F$.

\smallskip
\noindent\emph{Largeness.}
Let $Z\subset X$ be a positive-dimensional closed irreducible
subvariety.  If its projection to $B$ is nonconstant, then
$\pi_1(Z^\nu)\to\pi_1(B)$ has infinite image.  Otherwise its image,
a finite subgroup of the torsion-free group $\pi_1(B)$, would be
trivial.  The map $Z^\nu\to B$ would then lift holomorphically to
$\C^g$.  Every holomorphic function on the connected compact normal
space $Z^\nu$ is constant, a contradiction.

If the projection is constant, the fiber property gives $Z\subset F$.
For $Z=F$, its image in $\pi_1(X)$ is infinite cyclic by
\eqref{red:fundamental-group} and the hypothesis.  The only other
possibility is an irreducible curve in $A$.  Let $D$ be its smooth
normalization.  After a translation in $A$, the nonconstant map
$D\to A$ extends to a homomorphism $\Jac(D)\to A$.  Its image is a
nonzero abelian subvariety, and therefore all of $A$ by simplicity.
A surjective homomorphism of complex tori carries the source integral
lattice to a finite-index subgroup of the target lattice.  Consequently
$\pi_1(D)$ has finite-index image in $\pi_1(A)\simeq\Z^4$, and its
image under the homomorphism onto the infinite cyclic group remains
infinite.  This proves largeness in every case.  Lemma~\ref{red:compact-criterion}
then gives the required absence of compact analytic subvarieties in
$\widetilde X$.

\smallskip
\noindent\emph{The obstruction to Steinness.}
Let $K$ be the kernel of $\pi_1(A)\to\pi_1(X)$.  Its quotient is
infinite cyclic, so $K\simeq\Z^3$.  A connected component of the
inverse image of $F$ in $\widetilde X$ is the connected covering
\[
 A_K=\C^2/K.
\]
It is a closed complex submanifold of $\widetilde X$: the full inverse
image of $F$ is closed, and its components are closed and locally
separated in covering charts.  The real span of $K$ has dimension
three, so, as a real manifold,
\[
 A_K\simeq(\R^3/\Z^3)\times\R,
 \qquad H_3(A_K,\Z)\simeq\Z.
\]
By the Andreotti--Frankel theorem \cite{AndreottiFrankel1959}, a Stein
manifold of complex dimension two has the homotopy type of a CW complex of real dimension at most
two; see \cite[Theorem~7.2]{MilnorMorse}.  Thus $A_K$ is not Stein.
A closed complex submanifold of a Stein manifold is Stein, so
$\widetilde X$ cannot be Stein either.
\end{proof}

Proposition~\ref{prop:large-reduction} leaves one concrete task: embed a
simple abelian surface in a smooth projective variety so that its
fundamental group has infinite cyclic image.  We first construct the
corresponding map from a point blowup of the surface to a smooth
orbifold, and then convert it to the required embedding in
Section~\ref{sec:projective}.
